% Figura: problema 4.5. Cubo con polarización radial P = A r.
\begin{figure}[t]
\centering
\begin{tikzpicture}[scale=1.0, x={(1cm,0cm)}, y={(0.45cm,0.32cm)}, z={(0cm,1cm)}]
  \def\h{1.2}
  % caras traseras
  \draw[texto!45!fondo, line width=0.4pt, dashed] (-\h,\h,-\h) -- (\h,\h,-\h) (-\h,\h,-\h) -- (-\h,\h,\h) (-\h,\h,-\h) -- (-\h,-\h,-\h);
  % flechas radiales de P desde el centro
  \foreach \v in {{1,0,0},{-1,0,0},{0,0,1},{0,0,-1},{0.7,0,0.7},{-0.7,0,0.7},{0.7,0,-0.7},{-0.7,0,-0.7},{0,1,0},{0,-1,0}} {
    \draw[campoP, line width=0.8pt, -{Stealth[length=1.6mm]}] (0,0,0) -- ($0.95*\h*({\v})$);
  }
  % caras delanteras
  \fill[dielec, opacity=0.55] (-\h,-\h,-\h) -- (\h,-\h,-\h) -- (\h,-\h,\h) -- (-\h,-\h,\h) -- cycle;
  \fill[dielecB, opacity=0.5] (\h,-\h,-\h) -- (\h,\h,-\h) -- (\h,\h,\h) -- (\h,-\h,\h) -- cycle;
  \fill[dielecB, opacity=0.35] (-\h,-\h,\h) -- (\h,-\h,\h) -- (\h,\h,\h) -- (-\h,\h,\h) -- cycle;
  \draw[borde] (-\h,-\h,-\h) -- (\h,-\h,-\h) -- (\h,-\h,\h) -- (-\h,-\h,\h) -- cycle;
  \draw[borde] (\h,-\h,-\h) -- (\h,\h,-\h) -- (\h,\h,\h) -- (\h,-\h,\h);
  \draw[borde] (-\h,-\h,\h) -- (-\h,\h,\h) -- (\h,\h,\h);
  % normal en la cara x = l/2
  \draw[vecn] (\h,0,0) -- (\h+0.9,0,0) node[right, etiq] {$\nn=\uv x$};
  \node[etiq] at (0,-\h,-0.62*\h) {$\rhoP=-3A$};
  \node[etiq, anchor=west] at (\h+0.15,0,0.85*\h) {$\sigmaP=\tfrac12A\ell$};
  % ejes
  \draw[tenue, line width=0.4pt, -{Stealth[length=1.6mm]}] (\h,\h,-\h) ++(0.2,0,0) -- ++(0.7,0,0) node[right, etiqt] {$x$};
  \draw[tenue, line width=0.4pt, -{Stealth[length=1.6mm]}] (\h,\h,-\h) ++(0,0.3,0) -- ++(0,0.9,0) node[right, etiqt] {$y$};
  \draw[cota] (-\h,-\h,-\h-0.35) -- (\h,-\h,-\h-0.35) node[midway, below, etiq] {$\ell$};
\end{tikzpicture}
\caption{\Cref{prob:cubo}. Cubo de lado $\ell$ centrado en el origen con $\vP=A\vr$ ($A>0$). La polarización diverge
uniformemente, $\rhoP=-\divg\vP=-3A$, y en cada cara la componente normal de $\vP$ es constante: por ejemplo,
en $x=\ell/2$, $\sigmaP=\vP\cdot\uv x=A\ell/2$. La carga negativa del volumen compensa exactamente la positiva de las seis caras.}
\label{fig:p45}
\end{figure}
